\section{Sistema de variables y operacionalización}

\subsection{Definición conceptual de las variables}
% Defina cada variable con respaldo bibliográfico.
[Variable independiente: definición conceptual.]
[Variable dependiente: definición conceptual.]
[Variables secundarias: definición conceptual.]

\subsection{Operacionalización}
% La tabla debe permitir medir cada dimensión con indicadores concretos.
% Reemplace los textos entre corchetes.

\begin{table}[htbp]
  \centering
  \caption{Operacionalización de las variables}
  \label{tab:operacionalizacion}
  \small
  \renewcommand{\arraystretch}{1.15}
  \begin{tabularx}{\textwidth}{@{}>{\raggedright\arraybackslash}p{2.4cm}>{\raggedright\arraybackslash}p{3.2cm}>{\raggedright\arraybackslash}X>{\raggedright\arraybackslash}p{2.4cm}@{}}
    \toprule
    \textbf{Variable} & \textbf{Dimensión} & \textbf{Indicadores} & \textbf{Escala} \\
    \midrule
    [Variable independiente] & [Dimensión 1] & [Indicador 1; indicador 2] & [Nominal] \\
     & [Dimensión 2] & [Indicador 3] & [Ordinal] \\
    \midrule
    [Variable dependiente] & [Dimensión 1] & [Indicador 4; indicador 5] & [Razón] \\
     & [Dimensión 2] & [Indicador 6] & [Intervalo] \\
    \bottomrule
  \end{tabularx}
  \notaTabla{Elaboración propia a partir de [autor, año].}
\end{table}
